\chapter{複素関数}
    \section{Hermite対称関数と反Hermite対称関数への分解}
        \label{decomp_into_Hermite_sym_and_anti_Hermite_sym_fxns}
        \begin{shadebox}
            1変数関数$f$はHermite対称関数と反Hermite対称関数に1通りに分解できる
        \end{shadebox}
            \begin{proof}
                \quad\par
                (分解可能性)
                \par
                $\fs(x) \coloneq (f(x)+\conj{f(-x)})/2,\fa(x) \coloneq (f(x)-\conj{f(-x)})/2$とすれば確かに$f(x)=\fs(x)+\fa(x)$であり$\fs,\fs$はそれぞれHermite対称，反Hermite対称である。
                \newline
                (一意性)
                \par
                次に，分解が1通りであることを示す。
                Hermite対称関数$\fsI,\fsII$，反Hermite対称関数$\faI,\faII$が存在して$f(x) = \fsI(x) + \faI(x) = \fsII(x) + \faII(x)$と表せるとする。
                このとき$\fsI(x) - \fsII(x) = \faII(x) - \faI(x)$が成り立つ。
                左辺はHermite対称，右辺は反Hermite対称なので，両辺は恒等的に0でなければならない（Hermite対称かつ反Hermite対称な関数が恒等的に0であることは容易に示せる）。
            \end{proof}